% C.3af: full proof and compressed section construction independently audited.
\paragraph{C.3af (Refuted, with explicit counterexamples): ordinary restriction fails in every codimension in an explicit unbounded dimension range.}
Let $r\ge2$ and
\[
                         N\ge 4^{r-1}+2r-1.
\]
There is a completely specified excellent normal strictly
henselian local complex ring $R$ of dimension $N$, with
regular puncture $X$ and regular punctured principal divisor
$Y$, such that
\begin{equation}\label{eq:C3af-cokernel}
 \operatorname{coker}\!\left(
 \mathrm{CH}^r(X)_{\mathrm{alg}}
       \longrightarrow\mathrm{CH}^r(Y)_{\mathrm{alg}}\right)
          \longrightarrow\operatorname{Pic}^0(C)(\mathbb C)
\end{equation}
is surjective. Here $C$ is the Fermat plane quartic and
the equivalence relation is ordinary algebraic equivalence
through actual families on the fixed product. In particular,
the ordinary rational $J^r$ restriction is not surjective.
For $r=2$ this recovers the threshold $N\ge7$ of C.3ae.
No conclusion in the intermediate dimensions follows from
the bound in this statement.

\emph{1. An explicit full set of sections.}
We first use the larger bound
$N\ge3\cdot4^{r-1}+2r-1$ and then reduce it in step 5.
Set $s=r-1$, $q=3\cdot4^s$ and $m=N-2r\ge q-1$.
Let $W_1,\ldots,W_s$ be copies of the Fermat quartic K3
surface from C.3ae. On each $W_i$, write
\[
 H_i=c_1(\mathcal O_{W_i}(1)),\qquad
 D_i=V(w^{(i)}_0-\eta w^{(i)}_1,
        w^{(i)}_2-\eta w^{(i)}_3),\qquad
 \tau_i=4[D_i]-H_i,
 \quad\eta^4=-1.
\]
Put
\[
 \mathcal W=\prod_{i=1}^s W_i,\qquad
 S=\mathcal W\times C,\qquad
 M=\left(\boxtimes_{i=1}^s\mathcal O_{W_i}(1)\right)\boxtimes K_C,
 \qquad a=c_1(M).
\]
The $q$ sections
\begin{equation}\label{eq:C3af-sections}
       t_{\mathbf j,k}
          =z_k\prod_{i=1}^s w^{(i)}_{j_i},
 \qquad \mathbf j\in\{0,1,2,3\}^s,\quad k\in\{0,1,2\},
\end{equation}
form a basis of $H^0(S,M)$ and have no common zero.
Indeed at any point we choose one nonzero coordinate in
each projective factor and take their product. Give
$\mathbb P^m$ coordinates $x_{\mathbf j,k}$ for these
$q$ indices and, if needed, additional coordinates which
will not occur in
\[
                      \ell=\sum_{\mathbf j,k}
                            x_{\mathbf j,k}t_{\mathbf j,k}.
\]
Every coefficient is now zero or one; no generic choice
of a coefficient or a degree remains.

Take
\[
 B=\mathbb P^m\times\mathcal W\times\mathbb P^2,
 \qquad
 \mathcal N=\mathcal O_{\mathbb P^m}(-1)
       \boxtimes\left(\boxtimes_i\mathcal O_{W_i}(-1)\right)
       \boxtimes\mathcal O_{\mathbb P^2}(-1).
\]
In $\operatorname{Tot}(\mathcal N)$ define
\[
 Z=V(F+t\ell G),\qquad
 F=z_0^4+z_1^4+z_2^4,\qquad
 G=z_0^2z_1^2+z_0^2z_2^2+z_1^2z_2^2.
\]
Contract by all Segre coordinates
$t x_\alpha z_k\prod_i w^{(i)}_{j_i}$, take the reduced
closed image, normalize it, localize at the unique point
over the vertex, and take its strict henselization.
This defines $R$. The graph, Cohen--Macaulay,
normalization and scheme-theoretic exceptional-fibre
arguments in C.3ae apply verbatim to this explicitly
given smooth product. They give a projective regular
model $\mathcal Z$ with exceptional divisor
\[
 V=\mathbb P^m\times S,
 \qquad N_{V/\mathcal Z}=L^{-1},
 \qquad L=\mathcal O_{\mathbb P^m}(1)\boxtimes M.
\]
Its dimension is $m+2s+2=N$. In particular the
finite normalization has one vertex, and
$\mathfrak m_R\mathcal O_{\mathcal Z}=\mathcal I_V$;
the first neighbourhood is preserved by strict
henselization. The linear cone function $f=t\ell$ has
regular strict transform $\mathcal D$, with
\[
 H=V(\ell)\subset V,
 \qquad
 H=\mathbb P_{\mathrm{lines}}(E),
 \quad E=\ker(\mathcal O_S^{m+1}\to M),
 \quad\operatorname{rank}E=m.
\]
The map defining $E$ is surjective by
\eqref{eq:C3af-sections}. Its strict transform before
the local base change is the total space of $L_H^{-1}$.
Thus $X=\mathcal Z\setminus V$ and
$Y=\mathcal D\setminus H=X\cap V(f)$ have all the
asserted regularity properties.

\emph{2. The invariant in codimension $r$.}
The inequality $m\ge q-1>r$ implies that no
projective-bundle relation has degree at most $r$.
For every smooth projective parameter $P$, localization
and the refined self-intersection formula therefore give
\[
\begin{split}
 \sigma_P:\mathrm{CH}^r(Y\times P)&\longrightarrow
 \frac{\mathrm{CH}^r(H\times P)}
 {(h+a)\mathrm{CH}^{r-1}(H\times P)}\\
 &\simeq\mathrm{CH}^r(S\times P),\qquad h\longmapsto-a.
\end{split}
\]
The construction and independence are exactly those
proved in C.3ae, now with codimension $r$.
Intersect with $\tau=\prod_i\tau_i$ and push along
the $2s$-dimensional variety $\mathcal W$. Since
$r+s-2s=1$, this defines an integral homomorphism
\[
 \rho_P:\mathrm{CH}^r(Y\times P)
                \longrightarrow\operatorname{Pic}(C\times P)
\]
commuting with fibres. In particular $\rho$ maps the
ordinary algebraically trivial subgroup into
$\operatorname{Pic}^0(C)(\mathbb C)$.

For a closure of a cycle $\Gamma$ from $X\times P$,
write its refined restriction to $V\times P$ as
\[
 b=\sum_{j=0}^r h^j b_j,
 \qquad b_j\in\mathrm{CH}^{r-j}(S\times P).
\]
Transversality gives
\begin{equation}\label{eq:C3af-invariant}
 \rho_P(\Gamma|_{Y\times P})
       =\operatorname{pr}_{C\times P,*}
          \left(\tau\sum_{j=0}^r(-a)^j b_j\right).
\end{equation}
These refined pullbacks and pushforwards act on
actual Chow classes on the fixed product.

\emph{3. The only relevant Hodge component.}
All differential sheaves here are relative to
$\mathbb C$. The rank-one conormal sequence gives
\[
 0\longrightarrow L\otimes\Omega^{r-1}_{V\times P}
 \longrightarrow\Omega^r_{\mathcal Z\times P}|_{V\times P}
 \longrightarrow\Omega^r_{V\times P}\longrightarrow0.
\]
Its cohomology boundary is contraction with the
first-neighbourhood class $\kappa$. The rational
Chern-character lift from C.3ae applies in degree $r$:
the trace formula is
$\operatorname{ch}_r=\operatorname{tr}(\operatorname{At}^r)/r!$.
Consequently every closure restriction $b$ satisfies
\[
       \kappa\mathbin{\lrcorner}\operatorname{cl}_{\mathrm{Hdg}}(b)
                  =0\quad\hbox{in }
       H^{r+1}(V\times P,L\otimes\Omega^{r-1}_{V\times P}).
\]
Here the existence and pullback compatibility of the
virtual-bundle lift use the already opened
\citet[Lemma 42.58.1, Tag 0FEY;
Remark 42.59.9, Tag 0FF8]{cyStacksChernObstruction}.

The class $\kappa$ contracts only a $C$ differential;
it is
\[
          \kappa=\sum_{\mathbf j,k}x_{\mathbf j,k}
                   \left(\prod_i w^{(i)}_{j_i}\right)
                     \otimes\delta(z_kG).
\]
The three curve classes form the basis checked in
C.3ae. Project the vanishing to the Kunneth summand
\[
 H^0(\mathbb P^m,\mathcal O(1))\otimes
 \left(\bigotimes_{i=1}^sH^1(W_i,\Omega^1_{W_i}(1))\right)
 \otimes H^1(C,K_C)\otimes H^1(P,\mathcal O_P).
\]
Its cohomology degree is $s+1+1=r+1$ and its
differential degree is $s=r-1$. Since multiplication
by each $w^{(i)}_{j_i}$ preserves the cohomology and
differential degrees on $W_i$, the only source
component that can reach it is the component of $b_0$
in
\begin{equation}\label{eq:C3af-source}
 \left(\bigotimes_{i=1}^sH^1(W_i,\Omega^1_{W_i})\right)
         \otimes H^0(C,K_C)\otimes H^1(P,\mathcal O_P).
\end{equation}
All components involving a projective-space
differential are excluded. On a K3 factor the
$(2,0)$ and $(0,2)$ pieces have the wrong separate
degrees, and cannot enter this projection by
multiplication by a function. This excludes the
additional Kunneth pieces that appear on a product
of several K3 surfaces.

For each factor put
$U_i=H^1(W_i,\Omega^1_{W_i})$ and
$V_i=H^1(W_i,\Omega^1_{W_i}(1))$. C.3ae's Euler,
normal and Serre-duality calculation proves that
the joint multiplication map
\[
 U_i\longrightarrow\mathbb C^4\otimes V_i,
       \qquad u\longmapsto(w^{(i)}_0u,\ldots,w^{(i)}_3u)
\]
has kernel exactly $\mathbb C H_i$. Its quotient
is dual to the surjective multiplication
$\mathcal R_1\otimes\mathcal R_3\to\mathcal R_4$.
Surjectivity is elementary for the Fermat ring:
every degree-four monomial with all exponents at
most two is a variable times a surviving
degree-three monomial.

Because every section index in
\eqref{eq:C3af-sections} has its own independent
coordinate $x_{\mathbf j,k}$, the projected obstruction
is the tensor product of these joint multiplication
maps and the perfect curve pairing. Tensor products
over $\mathbb C$ preserve injections. Its kernel in
\eqref{eq:C3af-source} is therefore the sum of the
subspaces having $\mathbb C H_i$ in at least one
K3 factor. Pairing with $\tau_i$ in every factor
kills this entire kernel, since $\tau_iH_i=0$.
Thus the mixed $H^0(C,K_C)\otimes H^1(P,\mathcal O_P)$
component of $\operatorname{pr}_{C\times P,*}(\tau b_0)$
is zero.

We also check every correction term in
\eqref{eq:C3af-invariant}. The class $a$ is
$\sum_iH_i+c_1(K_C)$. To contribute to the specified
mixed Picard component after intersection with
$\prod_i\tau_i$, the class before that intersection
must have cohomology and differential degrees $(1,1)$
on every $W_i$, and degrees $(0,1)$ on $C$.
A factor $c_1(K_C)$ already has curve cohomology
degree one, so it cannot occur. A factor $H_i$
requires the rest of the tensor to have degrees
$(0,0)$ on that factor; its pairing is then
$\tau_iH_i=0$. Two such factors on one surface
have too large a degree. Thus every term with
$j\ge1$ contributes zero to the mixed component,
including terms whose $b_j$ contains nontrivial
cycles on several K3 factors. The line bundle
in \eqref{eq:C3af-invariant} has zero mixed
component.

\emph{4. Ordinary families and the target quotient.}
Take a generator of the ordinary algebraically
trivial source group, represented by a curve family.
Compactify that parameter smoothly, close the
cycle on the fixed product, and then on
$\mathcal Z$ times the proper curve. The invariant
defines a morphism from the parameter curve to
$\operatorname{Pic}^0(C)$ after normalization at
one point. Step 3 makes its induced homomorphism
from the parameter Jacobian zero on
$H^1(\mathcal O)$, hence zero. The family is
constant in $\operatorname{Pic}(C)$, proving
\[
             \rho\bigl(\mathrm{CH}^r(X)_{\mathrm{alg}}|_Y\bigr)=0.
\]

The product of the pulled-back divisors
$D_1,\ldots,D_s$ with a divisor class from
$\operatorname{Pic}^0(C)$ gives an actual ordinary
target family and a homomorphism
\[
 \theta:\operatorname{Pic}^0(C)(\mathbb C)
                    \longrightarrow\mathrm{CH}^r(Y)_{\mathrm{alg}}.
\]
As before, the Poincare parameter can be restricted
to chains of smooth connected curves. The exact
intersection identity is now
\[
                         \rho\theta=(-9)^s\,\mathrm{id}.
\]
Divisibility of the complex point group proves
the surjection \eqref{eq:C3af-cokernel} under the
temporary larger dimension bound. The same
specified point $p=(u:v:1)$, $v^4=-1-u^4$, with
$u$ transcendental over $\overline{\mathbb Q}$,
and $o=(1:\eta:0)$, gives a nontorsion target
cycle $\theta([p-o])$. This proves failure with
rational coefficients as well.

\emph{5. A smaller explicit set of sections.}
We now replace $3\cdot4^s$ by $4^s$ and retain the
same proof. Use the four matrices $A_i=(a_{ijk})$
from C.3ae, and define the $4$ by $3$ matrix
\[
                       B(w)_{ik}=\sum_{j=0}^3a_{ijk}w_j.
\]
At every point of the Fermat surface $W$, every
pair of columns of $B(w)$ is linearly independent.
This assertion has a separate exact certificate:
for each of the three column pairs and each of
the four charts $w_j=1$, the Fermat equation
together with the six $2$ by $2$ minors has
Groebner basis $[1]$ over $\mathbb F_{101}$.
There are twelve charts in total. Properness
of each projective determinantal locus over
$\mathbb Z$ proves the assertion in characteristic
zero, exactly as in C.3ae. The three coefficient
matrices for the individual columns have
determinants $-36$, $4$, and $200$, which also
verify that none of these columns vanishes at
a projective point. The retained computation
took $0.325551998$ seconds.

For $\mathbf i\in\{0,1,2,3\}^s$ set
\begin{equation}\label{eq:C3af-compressed-sections}
             t'_{\mathbf i}
                  =\sum_{k=0}^2 z_k
                        \prod_{l=1}^s B(w^{(l)})_{i_lk}.
\end{equation}
These $4^s$ sections of $M$ are base point free.
For $s=1$ this is the original base point
certificate in C.3ae. Suppose $s\ge2$. Write
$b_{lk}\in\mathbb C^4$ for the $k$th column
of $B(w^{(l)})$ at the chosen point. The three
vectors
\[
                         v_k=\bigotimes_{l=1}^s b_{lk}
                         \quad(k=0,1,2)
\]
are linearly independent. Here is a direct proof.
The three first factors $b_{1k}$ are pairwise
independent, and so are the three nonzero
remaining tensor factors
$\bigotimes_{l=2}^s b_{lk}$. Two pure tensors
cannot be proportional unless their nonzero
factors are proportional in every factor.
In a supposed dependence of the three $v_k$,
all coefficients must be nonzero, since each
pair is independent. Projecting the first
factor modulo the span of its first two
vectors shows that the third belongs to
that span. Both of its coefficients in that
span are nonzero. Comparing the first two
coefficients in the tensor product would
then make two remaining tensor factors
proportional, a contradiction. Thus the
$v_k$ are independent. Since the coordinate
vector $(z_0,z_1,z_2)$ is nonzero, its linear
combination of the $v_k$ is nonzero. These
are exactly the values of
\eqref{eq:C3af-compressed-sections}.

Put $q'=4^s$, take $m\ge q'-1$, and use
\[
                       \ell'=\sum_{\mathbf i}
                                      x_{\mathbf i}t'_{\mathbf i}
\]
instead of $\ell$ in the geometric construction.
Its coefficients are the specified products
$\prod_l a_{i_lj_lk}$, so the construction is
still completely determined by the four
integer matrices. Base point freeness gives
the same regular divisor model, and
$m\ge4^{r-1}-1>r$ for $r\ge2$ gives the
same projective-bundle quotient.

It remains to check the obstruction, since
the new sections have correlated coefficients.
Let $\overline U=U/\mathbb C H_W$ and let
$E_C$ be the three-dimensional space with
basis $e_k$ dual to $\delta(z_kG)$. The
rank-$57$ certificate of C.3ae says precisely
that
\[
 \iota:\overline U\otimes E_C
       \longrightarrow\mathbb C^4\otimes H^1(W,\Omega_W^1(1)),
 \qquad
 u\otimes e_k\longmapsto
       \left(\sum_j a_{ijk}w_ju\right)_{i=0}^3
\]
is injective. For the new sections, the map
on the primitive quotient of
\eqref{eq:C3af-source} is the composition,
up to reordering tensor factors, of
\[
\begin{split}
 \overline U^{\otimes s}\otimes E_C
   &\longrightarrow(\overline U\otimes E_C)^{\otimes s}\\
   &\xrightarrow{\ \iota^{\otimes s}\ }
       \left(\mathbb C^4\otimes
              H^1(W,\Omega_W^1(1))\right)^{\otimes s}.
\end{split}
\]
The first map sends
$u_1\otimes\cdots\otimes u_s\otimes e_k$
to $\bigotimes_l(u_l\otimes e_k)$. It is
injective because the three tensors
$e_k^{\otimes s}$ are linearly independent.
The second is injective because tensor
products over $\mathbb C$ preserve
injections. Expanding the composition gives
exactly the coefficients
$\prod_l a_{i_lj_lk}$ in
\eqref{eq:C3af-compressed-sections}; there is
no replacement of the actual curve by a
product of curves. Therefore the only
kernel before the quotient still consists
of tensors with at least one $H_i$ factor.
The Hodge projection, all boundary
corrections and the Picard-family argument
of steps 3--4 apply unchanged. This proves
\eqref{eq:C3af-cokernel} for the asserted
smaller bound $N\ge4^{r-1}+2r-1$.

The proof uses the finite Jacobian-ring spanning
certificate in C.3ae and the tensor-product
argument for arbitrary $r$. It does not rely on
testing a finite list of codimensions or on
unbounded matrix computations. The coefficient
choice and the dimension bound remain explicit
for every $r$.
